\begin{alist}
\item 
\begin{rlist}
\item Γνωρίζουμε ότι για κάθε $x > 0$, $\alpha > 0$ και $\alpha \neq 1$ ισχύει ότι $\log_\alpha \alpha^x = x$. \\
Οπότε για $x = \alpha = 10$ έχουμε ότι $\log 10^{10} = 10$. \\
Εναλλακτικά, $\log 10^{10} = 10 \cdot \log 10 = 10 \cdot 1 = 10$.

\item Είναι $A = \ln(\ln e) + \log(\log 10^{10}) = \ln 1 + \log 10 = 0 + 1 = 1$.
\end{rlist}

\item Η εξίσωση $\log(x^2 + 1) = A$ ορίζεται για κάθε $x \in \mathbb{R}$, αφού $x^2 \ge 0 \Leftrightarrow x^2 + 1 \ge 1 > 0$ για κάθε $x \in \mathbb{R}$ και ισοδύναμα έχουμε:
\[\log(x^2 + 1) = A \Leftrightarrow \log(x^2 + 1) = 1 \Leftrightarrow x^2 + 1 = 10 \Leftrightarrow x^2 = 9 \Leftrightarrow x = \pm 3.\]
\end{alist}
